\section{Introduction}

Collective motion is often studied through systems that display visible global
order, such as polarized schools, flocks, or moving groups. Laboratory midge
swarms provide a different test case. The insects move in a disordered
three-dimensional cloud and do not exhibit a strong global orientational order,
yet the swarm remains spatially cohesive and shows reproducible internal
structure. This combination raises a basic question: at what level should the
organization of such a weakly ordered group be described?

The question is not trivial because position-dependent velocity differences
are expected even in the absence of a new behavioral mechanism. Translation of
the swarm center, local rotation, expansion, compression, shear, and density
gradients can all make insects at different positions move differently. Thus,
observing that local velocities differ across the swarm is not by itself a
strong result. A more informative test is to ask whether any stable component
of motion remains after ordinary local kinematic explanations have been given
a clear baseline.

\subsection{Existing Descriptions of Midge Swarms}

Previous work has shown that laboratory midge swarms are not uniform random
particle clouds. Time-resolved three-dimensional trajectory data provide the
basis for measuring individual motion, local neighborhoods, and global swarm
geometry \cite{sinhuber2019_dataset}. Individuals also sample the swarm volume
non-uniformly, so spatial position and local density must be treated as
potential contributors to any local motion pattern
\cite{feng2023_sampling}.

A second line of work has described swarms using macroscopic and
statistical-physics variables. Swarm size, radius, density, compactness, and
layer structure provide compact summaries of the full trajectory data, and
equation-of-state and thermodynamic descriptions show that weakly ordered
swarms can nevertheless possess reproducible collective states
\cite{sinhuber2021_eos,reynolds2021_thermodynamic}. These studies make
low-dimensional state variables natural explanatory candidates for local
motion, but they do not by themselves determine whether a local residual can
be accounted for by those variables.

A third line of work treats midge motion in stochastic dynamical terms.
Restoring-force descriptions, Langevin-type models, and intrinsic
stochasticity have been used to connect microscopic fluctuations with
macroscopic swarm organization
\cite{reynolds2017_velocity_gravity,reynolds2018_langevin,reynolds2021_intrinsic}.
This perspective is important for the present study because a residual signal
need not imply a deterministic rule. It may reflect incomplete geometric
subtraction, a compact stochastic state description, or local organization that
is not captured by the tested variables.

Finally, perturbation and correlation analyses show that midge swarms contain
measurable spatiotemporal structure. Environmental perturbations can induce
correlated responses, and spatial correlations have been reported in
laboratory insect swarms
\cite{vandervaart2020_perturbations,reynolds2024_spatial}. Such correlations
motivate a local analysis, but they do not distinguish local coordination from
ordinary geometric deformation, sampling structure, or density-dependent
motion. This leaves open the specific residual question addressed here.

\subsection{Non-Affine Residuals as a Methodological Reference Point}

A useful methodological precedent comes from non-affine analysis in complex
many-body systems. In amorphous solids and colloidal glasses, local affine
motion is often fitted first, and the remaining non-affine displacement is then
used to characterize local rearrangement
\cite{falk1998_viscoplastic,chikkadi2012_nonaffine}. Related decompositions
have also been adapted to biological collective motion, including epithelial
sheet migration and cell assemblies, where large-scale collective motion must
be separated from local rearrangements
\cite{lee2013_epithelial,cerbino2021_disentangling}.

The role of this literature in the present paper is methodological rather than
ontological. We do not treat a midge swarm as a glass, granular material, or
elastic continuum. Instead, we borrow a more general diagnostic logic: first
allow a local neighborhood to explain as much relative motion as possible
through a smooth affine deformation field, and then analyze only the part that remains. This
choice makes the baseline conservative. Local rotation, extension,
compression, and shear are not immediately interpreted as new collective
organization; they are first absorbed by the affine fit.

\subsection{From Candidate Explanations to T1}

The present analysis was motivated by a sequence of candidate explanations.
First, we asked whether the transition-linked tangential signal could be
absorbed by ordinary geometry, including global motion and local affine
deformation. If local affine geometry removed the signal, no additional local
residual object would be needed. Second, when a residual remained, we asked
whether it was stable to nearby choices of local scale and temporal lag. Third,
we examined whether the residual had a repeated spatial or temporal form, such
as an edge-localized contrast, a sharp event precursor, or a signed transition
pattern. Finally, we tested whether a small set of state variables, true event
timing, or recent state-path history was sufficient to account for it.

This paper therefore does not start from a proposed complete mechanism. It
starts from a residual object that survives several natural explanations. For
a focal individual \(i\) at time \(t\), we describe each retained neighbor
\(j\in N_i(t)\) by the focal-centered relative vector
\(\mathbf{r}_{ij}(t)=\mathbf{x}_j(t)-\mathbf{x}_i(t)\). Over a finite lag
\(\ell\), the neighborhood deformation is summarized by
\(\Delta\mathbf{r}_{ij}(t,\ell)=\mathbf{r}_{ij}(t+\ell)-\mathbf{r}_{ij}(t)\).
The local affine fit is written conceptually as
\begin{equation}
\hat{\mathbf{J}}_i(t)
=
\arg\min_{\mathbf{J}}
\sum_{j\in N_i(t)}
\left\|
\Delta\mathbf{r}_{ij}(t,\ell)
-
\mathbf{J}\mathbf{r}_{ij}(t)
\right\|^2 .
\end{equation}
The non-affine residual is the part of neighbor relative motion not explained
by this local affine deformation,
\begin{equation}
\mathbf{u}^{NA}_{ij}(t)=
\frac{
\Delta\mathbf{r}_{ij}(t,\ell)
-
\hat{\mathbf{J}}_i(t)\mathbf{r}_{ij}(t)
}{\Delta t_{\ell}}.
\end{equation}
T1 denotes the tangential activity of this focal-neighborhood residual around
compact-density transition windows. When the analysis uses unsigned activity,
T1 is summarized by the magnitude or squared magnitude of the tangential
residual. The essential point is that T1 is not a raw speed, an inferred force,
or an individual behavioral rule. It is a local tangential non-affine residual
measured after local affine deformation has been removed.

\subsection{Aim and Contribution}

This study asks four questions. First, does T1 remain detectable in most
observations after local affine subtraction? Second, is that survival stable to
nearby choices of local scale and temporal lag? Third, what spatial and
temporal form does the surviving residual take? Fourth, are compact state
variables, true transition timing, or recent path history sufficient to
account for it?

The contribution is deliberately bounded. We do not claim to identify a
universal mechanism for midge swarming, nor do we reject existing stochastic or
statistical-physics descriptions of these systems. Instead, we define and test
a local tangential non-affine residual across 19 laboratory midge-swarm
observations. T1 survives local affine subtraction in most observations and is
robust within the survivor class, with diffuse tangential activity as the most
stable repeated form. At the same time, the tested compact state-variable,
event-timing, and recent-history explanations do not account for it
consistently across observations. The resulting picture is a reproducible
local collective observable with explicit observation-level boundaries, rather
than a completed mechanism model.
